% ============================================================
% CAPITOLO 17 - RIEPILOGO DELLA PARTE IV
% ============================================================
\chapter{Riepilogo della Parte IV}
\label{chap:riepilogo_iv}

\section{Quadro delle progressioni di triadi}
\label{sec:quadro_progressioni_triadi}

\begin{center}
\small
\begin{tabular}{p{5.2cm}p{6.8cm}}
\toprule
\textbf{Progressione} & \textbf{Dove} \\
\midrule
Definizione, elementi, errore delle ripetizioni eccessive & Sezione \ref{sec:definizione_progressione} \\
Quinte discendenti & Sezione \ref{subsec:progressione_quinte} \\
Terze discendenti & Sezione \ref{subsec:progressione_terze_disc} \\
Moduli I--IV, ascendente e discendente & Sezione \ref{sec:moduli_i4} \\
Moduli I--V, ascendente e discendente & Sezione \ref{sec:moduli_i5} \\
Progressioni imitate & Sezione \ref{sec:progressioni_imitate} \\
\bottomrule
\end{tabular}
\end{center}

Le progressioni che usano le settime sono nel Capitolo \ref{chap:progressioni_settime}, quelle che modulano nel Capitolo \ref{chap:progressioni_modulanti}, quelle cromatiche nel Capitolo \ref{chap:progressioni_cromatiche}. Il Capitolo \ref{chap:riepilogo_viii} riunisce tutte le progressioni in un'unica guida alla scelta.
